%% kapa: exam=1 round=30 year=2019 period=1 form=A subject=economics subject_ko=경제학원론 seq=16 number=56
%% source: Q-Net 감정평가사 기출문제 (cst00309) artlSeq=5205205 file=extracted/1008/5205205/2019년 제30회 감정평가사 1차시험 문제지/2019년 제30회 감정평가사 1차시험 문제지/2019년도 제30회 감정평가사 1차 1교시 A형.hwp
%% conversion: kordoc 4.24.51 -> tools/gen1.py (auto); review=proofread-sonnet; flags=none
\begin{question}{56}
다음의 전략형 게임(strategic form game)에서 $\alpha$에 따라 甲과 乙의 전략 및 균형이 달라진다. 이에 관한 설명으로 옳지 않은 것은? (단, 보수 행렬의 괄호 안 첫 번째 보수는 甲, 두 번째 보수는 乙의 것이다.) \par
{\small\centering
\begin{tabular}{|c|c|c|c|}
\hline
\multicolumn{2}{|c|}{\multirow{2}{*}{}} & \multicolumn{2}{|c|}{\makecell{乙}} \\
\cline{3-4}
\multicolumn{2}{|c|}{} & \makecell{Left} & \makecell{Right} \\
\hline
\multirow{2}{*}{\makecell{甲}} & \makecell{Up} & \makecell{(5-$\alpha$, 1)} & \makecell{(2, 2)} \\
\cline{2-4}
 & \makecell{Down} & \makecell{(3, 3)} & \makecell{(1, $\alpha$-1)} \\
\hline
\end{tabular}\par}

\begin{choices}
  \item $\alpha <2$이면, 전략 Up은 甲의 우월전략이다.
  \item $\alpha >4$이면, 전략 Right는 乙의 우월전략이다.
  \item $2< \alpha <4$이면, (Down, Left)는 유일한 내쉬균형이다.
  \item $\alpha <2$이면, (Up, Right)는 유일한 내쉬균형이다.
  \item $\alpha >4$이면, (Up, Right)는 유일한 내쉬균형이다.
\end{choices}
\end{question}
\answer{3}
